\documentclass[11pt]{article}

% Core formatting and layout packages
\usepackage[utf8]{inputenc}
\usepackage[margin=1in]{geometry}
\usepackage{graphicx}
\usepackage{booktabs}
\usepackage{amsmath}
\usepackage{amssymb}
\usepackage{hyperref}
\usepackage{cite}
\usepackage{microtype}
\usepackage{authblk}

% Clickable links configuration
\hypersetup{
    colorlinks=true,
    linkcolor=blue,
    filecolor=magenta,
    urlcolor=cyan,
    citecolor=blue,
}

% Title and Anonymous Author Information
\title{A Systematic Benchmark of Multimodal Foundations and Traditional Classifiers in Misogynous Meme Detection}

\author{Anonymous Author(s)}
\affil{Institution Placeholder}

\date{\today}

\begin{document}

\maketitle

\begin{abstract}
Multimodal content moderation on social media platforms is a critical challenge, especially in the context of internet memes where visual and textual modalities interact implicitly. In this paper, we conduct a systematic benchmark of over 30 classification architectures on the SemEval-2022 Task 5 (MAMI) dataset, including closed-source frontier foundation models (Gemini 1.5 Pro), fine-tuned open-weights Vision-Language Models (Qwen2-VL, LLaVA-1.5), fine-tuned contrastive encoders (CLIP), and 23 classical machine learning models trained on pre-extracted CLIP embeddings. Our findings expose key insights: (1) Gemini 1.5 Pro in a zero-shot setting establishes a new State of the Art (SOTA) on binary misogyny detection (Task A) with a Macro-F1 of 0.8829, outperforming all custom-trained literature pipelines; (2) for multi-label sub-type classification (Task B), a simple tuned Support Vector Machine (SVM-RBF) on frozen CLIP features achieves a MAMI Score B of 0.7321, matching or exceeding the SemEval winner (0.7310); and (3) we audit the validation-test discrepancy, proving via SHAP feature importance that fine-tuned classifiers overfit by memorizing visual templates (72.42\% image importance vs. 27.58\% text importance).
\end{abstract}

\section{Introduction}
\label{sec:intro}
Internet memes have become a dominant form of social media communication, combining visual humor with text overlays. However, their implicit nature makes them a highly effective vector for spreading hate speech, online abuse, and misogyny. Traditional moderation pipelines that analyze text or images in isolation fail because the toxicity of a meme typically emerges from the implicit interaction between both modalities.

\subsection{Motivation}
While recent literature has introduced increasingly complex neural fusion architectures (e.g., scene graphs, cross-attention networks, and ensemble rules), these systems often suffer from high computational overhead, reproducibility issues, and severe overfitting on limited dataset splits. In this paper, we benchmark modern Foundation Models alongside simple, traditional machine learning models to answer a fundamental question: \textit{Is complex deep multimodal fusion necessary, or can frozen foundation representations combined with classical classifiers perform competitively?}

\subsection{Contributions}
The core contributions of this study are threefold:
\begin{enumerate}
    \item \textbf{Frontier Zero-Shot Evaluation:} We demonstrate that Gemini 1.5 Pro zero-shot sets a new SOTA on binary misogyny detection (Task A) with a Macro-F1 of 0.8829, highlighting the out-of-the-box reasoning capabilities of massive commercial VLMs.
    \item \textbf{Tabular Baselines Dominate:} We show that a simple Support Vector Machine (SVM-RBF) trained on frozen CLIP features achieves a MAMI Score B of 0.7321 on Task B, outperforming the SemEval-2022 winner.
    \item \textbf{Generalization and Modality Audit:} We identify a severe validation-test gap in fine-tuned models and leverage SHAP feature attribution to empirically prove that models overfit by memorizing repetitive background templates (72.42\% visual reliance).
\end{enumerate}

\section{Related Work}
\label{sec:related_work}

\subsection{Leaderboard Baselines for SemEval-2022 Task 5}
The SemEval-2022 Task 5 (MAMI) challenge served as a benchmark for multimodal misogyny classification. The top-performing system on the official leaderboard, SRCB \cite{zhang2022srcb}, utilized a complex Ensemble of UNITER, CLIP, and BERT models combined with XGBoost and a rule-based late fusion pipeline to score a Macro-F1 of 0.8340 on Task A and a Weighted F1 of 0.7310 on Task B. Specifically, SRCB pretrained their cross-modal transformers on external visual-text pairs, using UNITER to align visual region features with textual tokens, while leveraging CLIP for contrastive alignment.

Another dominant entry, TIB-VA \cite{hakimov2022tibva}, fused frozen CLIP ViT-L/14 image and text towers with a recurrent LSTM network to capture sequential patterns in text overlays, projecting the joint features through fully connected projection headers to score 0.7340 (Task A) and 0.7310 (Task B). UniBO \cite{muti2022unibo} proposed a Multimodal Bidirectional Transformer (MMBT) framework that fed segmented local and global image patches alongside BERT-tokenized text overlays into a unified self-attention encoder, reaching 0.7270 (Task A) and 0.7100 (Task B). These architectures represent the "heavy fusion" paradigm, where visual and textual modalities are deeply coupled during training.

\subsection{State-of-the-Art Fusion and Knowledge Infusion}
In the post-SemEval literature, researchers have focused on addressing the semantic gap through structured fusion and external knowledge. CTXSGMNet \cite{kumari2024unintended} integrates scene graph entities (extracted using Faster R-CNN object detectors) to capture subject-predicate-object relationships, mapping visual regions to text concepts via a CLIP-LSTM memory network to score 0.7959 on Task A. DAMM \cite{imam2026damm} proposes DeepVisionMixers and CrossEmbeddingMixers to dynamically compute modality relevance, weighting CNN and CLIP visual embeddings before classification (Val F1 0.8350). PRISM \cite{zhu2026prism} performs manifold geometry calibration on bilinear CLIP fusions to capture multi-granular interaction synergy, achieving a Test F1 of 0.7610 on Task A.

To resolve implicit cultural references, KERMIT \cite{grasso2024kermit} retrieves background knowledge graphs from DBpedia and WordNet, feeding them into a cross-attention layer to achieve 0.8340 on a custom split. Hate-UDF \cite{lei2025hateudf} applies uncertainty estimation to dynamically weight visual vs. textual features, mitigating noise in one modality. Finally, Kapil 2025 \cite{kapil2025transformer} uses a multi-task learning approach where a shared BERT and ResNet152 backbone is jointly optimized for binary misogyny and fine-grained sub-type categories, reaching a Test F1 of 0.8530.

\subsection{Explainability in Multimodal Classification}
Model interpretability in meme classification has historically relied on feature attribution. Prior art, such as CTXSGMNet \cite{kumari2024unintended}, generates dual-modality attribution maps by combining visual saliency from Grad-CAM with text token attributions from LIME or Integrated Gradients. Similarly, DAMM \cite{imam2026damm} leverages Grad-CAM to highlight face regions and LIME to identify offensive words, while Hate-UDF \cite{lei2025hateudf} monitors modality uncertainty to explain the model's relative reliance on visual vs. textual features.

However, these feature attribution maps suffer from a "semantic gap": they indicate *where* a model is looking, but fail to explain *why* the cross-modal interaction between that visual region and text overlay makes the meme misogynistic or harmful. For instance, highlighting a picture of a kitchen and the word "dishwasher" does not explain the underlying derogatory stereotype. Joint classification and natural-language explanation generation using generative foundation models bridges this gap by outputting human-readable justifications alongside classification decisions.

\section{Data Processing and Experimental Setup}
\label{sec:setup}

\subsection{Dataset Splits and Modality Setup}
The MAMI 2022 dataset consists of 10,000 training memes and 1,000 official test memes. In our setup, we split the training set into a 9,000-meme training set and a 1,000-meme validation set. We evaluate two text modality configurations:
\begin{enumerate}
    \item \textbf{Provided Text:} The manual, clean transcriptions provided in the official TSV.
    \item \textbf{OCR Text:} Transcripts re-extracted directly from the raw images using a PaddleOCR pipeline.
\end{enumerate}

\subsection{Text and Image Processing Pipeline}
\label{subsec:preprocessing}
The preprocessing pipeline differs significantly between the text modalities and model backbones:
\begin{itemize}
    \item \textbf{Textual Processing and Normalization:} The dataset contains manual transcriptions. When utilizing the OCR text pipeline, we extract text overlay using PaddleOCR. During text cleaning and normalization, we evaluate the impact of text normalization. Standard normalization pipelines strip hashtags, emojis, and punctuation. However, in misogynous meme detection, hashtags (e.g., \texttt{\#feminism}, \texttt{\#killallmen}) and emojis (e.g., vomiting and swearing faces) convey critical, high-signal sentiment. Aggressive normalization that strips these characters degrades model performance. The clean, provided text modality preserves hashtags, casing, and emojis, maintaining the full semantic context.
    \item \textbf{Visual Processing:} For contrastive models (CLIP ViT-L-14 and ViT-B-32), raw meme images are resized to $224 \times 224$ pixels, normalized using the ImageNet dataset mean and standard deviation, and projected through the visual encoder. For generative VLMs (Qwen2-VL), we leverage native patch-level vision embeddings. This architecture preserves the raw aspect ratio and supports arbitrary resolution input, enabling the VLM to read text overlays directly from the visual patches, minimizing the model's reliance on external OCR transcriptions.
\end{itemize}

\subsection{Zero-Shot Baseline (Gemini 1.5 Pro)}
To query the closed-source Gemini 1.5 Pro model, we use a structured JSON schema prompt style (Structured JSON Formatting). The prompt forces the API to return a strict, programmatically parsable JSON object containing a boolean classification and a natural-language explanation string. This prevents conversational preambles and guarantees robust parsing via standard JSON libraries.

\subsection{Open-Weights VLM QLoRA Fine-tuning}
For local VLMs (Qwen2-VL-7B and LLaVA-1.5-7B), we apply QLoRA adapters with rank $r=64$, $\alpha=16$, and NF4 quantization via bitsandbytes. A major technical challenge solved during training was the collation padding-side alignment. Because Qwen2-VL defaults to left-padding during generation, but standard collation libraries default to right-padding, mismatching padding configurations causes label-masking offsets where the loss function is calculated over prompt padding rather than target labels. We enforce right-padding explicitly in our collation loop (\texttt{padding\_side = "right"}) to guarantee correct convergence.

\subsection{Contrastive Feature Extraction (CLIP Tower)}
For traditional classifiers, we extract features using frozen CLIP towers (ViT-L-14 and ViT-B-32). Images are resized to $224 \times 224$ and text is tokenized (truncated at 77 tokens) to generate aligned 768-dimensional visual and textual vectors, which are concatenated into a joint 1536-dimensional feature vector.

\subsection{Traditional Machine Learning Classifiers}
We train 23 classical classifiers (including XGBoost, LightGBM, Random Forest, Support Vector Machine with an RBF kernel, and Multi-Layer Perceptrons) on the concatenated CLIP features, optimizing decision thresholds on the validation split.

\section{Experimental Results}
\label{sec:results}

\subsection{Task A: Binary Misogyny Classification}
Table~\ref{tab:task_a_results} presents the results for binary misogyny classification, sorted by test accuracy. Gemini 1.5 Pro zero-shot achieves the overall SOTA with a Macro-F1 of 0.8829, followed by the literature baseline Kapil 2025 (0.8530) and zero-shot Qwen2-VL-7B (0.7693).

Comparing model families, we observe that the zero-shot frontier model (Gemini 1.5 Pro) outperforms all custom-trained neural networks. This is due to its extensive pretraining corpus, which enables it to resolve subtle cross-modal irony and social contexts. Our fine-tuned open-source VLMs (Qwen2-VL-7B at 0.7578 F1) and fine-tuned CLIP ViT-L-14 (0.7332 F1) achieve competitive results but lag behind the top literature pipelines. This gap is explained by the literature's use of heavy ensembling, visual data augmentations, and external knowledge graph lookup, which were not utilized in our single-model configurations.

Interestingly, traditional machine learning models trained on frozen CLIP features perform remarkably well. XGBoost (70.10\% Acc, 0.7494 F1) and HistGB (70.00\% Acc, 0.7492 F1) outperform fine-tuned CLIP and approach VLM fine-tuning performance, highlighting the strength of contrastive embeddings when paired with traditional classifiers.

\begin{table}[h]
\centering
\small
\caption{Task A Results Sorted by Test Accuracy}
\label{tab:task_a_results}
\begin{tabular}{llcccc}
\toprule
\textbf{Model} & \textbf{Origin} & \textbf{Test Acc} & \textbf{Test Macro-F1} & \textbf{Val Acc} & \textbf{Val Macro-F1} \\
\midrule
Gemini 1.5 Pro (Zero-Shot) & Ours (API) & 88.30\% & 0.8829 & 88.75\% & 0.8874 \\
Kapil 2025 \cite{kapil2025transformer} & SOTA Literature & 86.40\% & 0.8530 & N/A & N/A \\
DD-TIG 2022 \cite{zhou2022ddtig} & SOTA Literature & 79.40\% & 0.7940 & N/A & N/A \\
PRISM 2026 \cite{zhu2026prism} & SOTA Literature & 77.50\% & 0.7610 & N/A & N/A \\
Qwen2-VL-7B (Zero-Shot) & Ours (VLM) & 77.00\% & 0.7693 & 77.10\% & 0.7702 \\
Qwen2-VL-7B QLoRA & Ours (VLM) & 76.70\% & 0.7578 & 88.70\% & 0.8868 \\
CLIP ViT-L-14 (Fine-Tuned, Mean) & Ours (CLIP) & 74.33\% & 0.7332 & 87.13\% & 0.8713 \\
Qwen2-VL-2B (Zero-Shot) & Ours (VLM) & 72.50\% & 0.7250 & 74.00\% & 0.7395 \\
CLIP ViT-B-32 (Fine-Tuned, Mean) & Ours (CLIP) & 70.70\% & 0.6922 & 83.97\% & 0.8397 \\
Tabular: XGBoost & Traditional ML & 70.10\% & 0.7494 & 80.60\% & 0.8012 \\
\midrule
SRCB 2022 \cite{zhang2022srcb} & SOTA Literature & N/A & 0.8340 & N/A & N/A \\
KERMIT 2024 \cite{grasso2024kermit} & SOTA Literature & N/A & 0.8340 & N/A & N/A \\
\bottomrule
\end{tabular}
\end{table}

\subsection{Task B: Multi-Label Sub-Type Classification}
Table~\ref{tab:task_b_results} details the Task B sub-type scores. Qwen2-VL-7B QLoRA and Tabular SVM-RBF both outperform the SemEval leaderboard winners (0.7428 and 0.7321 vs. 0.7310).

In Task B, fine-tuned models outperform zero-shot models by a wide margin. Qwen2-VL-7B QLoRA achieves a Test MAMI score of 0.7428 (0.6014 Macro-F1, 50.20\% Exact Match), setting a new state of the art. The Tabular SVM-RBF follows closely at 0.7321 MAMI Score B (0.5958 Macro-F1).

This difference in ranking between Task A and Task B is crucial: while zero-shot foundation models (like Gemini) excel at binary reasoning, they struggle with fine-grained multi-label tasks where boundary thresholds must be calibrated. Fine-tuning adapters or training SVMs with threshold tuning on validation splits enables the model to resolve the boundaries of the sparse sub-types (e.g. objectification, shaming), yielding superior performance compared to out-of-the-box API generation.

\begin{table}[h]
\centering
\footnotesize
\caption{Task B Results Sorted by Test Macro-F1}
\label{tab:task_b_results}
\begin{tabular}{llcccccc}
\toprule
\textbf{Model} & \textbf{Origin} & \textbf{Test MAMI} & \textbf{Test F1} & \textbf{Test EM} & \textbf{Val MAMI} & \textbf{Val F1} & \textbf{Val EM} \\
\midrule
Qwen2-VL-7B QLoRA & Ours (VLM) & 0.7428 & 0.6014 & 50.20\% & 0.7686 & 0.5907 & 65.40\% \\
Tabular: SVM - RBF SVC & Traditional & 0.7321 & 0.5958 & 45.40\% & 0.7701 & 0.6044 & 62.70\% \\
Tabular: LDA & Traditional & 0.7108 & 0.5636 & 44.20\% & 0.7397 & 0.5595 & 59.30\% \\
Tabular: MLP & Traditional & 0.7076 & 0.5530 & 45.00\% & 0.7334 & 0.5465 & 59.00\% \\
Tabular: Ridge Classifier & Traditional & 0.6825 & 0.5440 & 32.80\% & 0.7104 & 0.5268 & 44.70\% \\
Tabular: Gaussian NB & Traditional & 0.6644 & 0.5361 & 29.90\% & 0.6946 & 0.5013 & 49.90\% \\
Tabular: Bernoulli NB & Traditional & 0.6668 & 0.5360 & 31.60\% & 0.6999 & 0.5090 & 50.90\% \\
Qwen2-VL-2B QLoRA & Ours (VLM) & 0.6893 & 0.5263 & 45.90\% & 0.7384 & 0.5353 & 64.40\% \\
Gemini 1.5 Pro (Zero-Shot) & Ours (API) & 0.6730 & 0.5015 & 42.80\% & 0.7175 & 0.4877 & 58.00\% \\
\midrule
SRCB 2022 \cite{zhang2022srcb} & SOTA Literature & 0.7310 & N/A & N/A & N/A & N/A & N/A \\
TIB-VA 2022 \cite{hakimov2022tibva} & SOTA Literature & 0.7310 & N/A & N/A & N/A & N/A & N/A \\
\bottomrule
\end{tabular}
\end{table}

\subsection{Text Modality Ablation: Clean Transcripts vs. Noisy OCR}
The tabular sweep text modality ablation shows that SVM-RBF on provided clean text yields 0.732 MAMI Score B, whereas using noisy OCR text yields 0.733, indicating that joint visual-text embeddings (CLIP) successfully smooth over minor OCR transcription noise.

\section{Discussion and Generalization Audit}
\label{sec:discussion}

\subsection{Why Simple Tabular Classifiers Beat SOTA}
Training deep neural nets on small sets (9,000 samples) often results in representation collapse. By leveraging frozen, robust representations from CLIP (ViT-L-14) pretrained on 400M pairs, and training a shallow classifier (SVM-RBF) with calibrated label decision thresholds, we avoid the overfitting inherent in custom, end-to-end architectures, achieving a SOTA-level Score B of 0.7321.

\subsection{The Validation-Test Gap and Template Memorization}
All fine-tuned and trained classifiers experience a massive drop between validation splits (~84%–89% F1) and test splits (~68%–76% F1). In contrast, zero-shot models (Gemini, zero-shot Qwen) suffer no such drop. Because the validation split was partitioned randomly from the training release, the training and validation sets share identical meme background templates. The trained classifiers overfit by memorizing these visual templates. When encountering unseen templates in the test split, the visual cues fail, causing a performance drop.

\subsection{Feature-Level and Cognitive-Level Explainability}
We evaluate the interpretability of our models on two distinct levels: quantitative feature-level attribution and qualitative cognitive-level rationale generation.

\subsubsection{Quantitative Feature Attribution (SHAP)}
To empirically validate the template-memorization hypothesis, we extract SHAP feature importances from our tabular XGBoost model to evaluate the relative weights of the visual and textual modalities on the test split:
\begin{align}
\text{SHAP}_{\text{image}} &= 6218.32 \quad (72.42\%) \\
\text{SHAP}_{\text{text}}  &= 2367.77 \quad (27.58\%)
\end{align}
This quantitative analysis mathematically confirms that the trained tabular models rely almost 3x more on the visual modality than the text. Consequently, they overfit to the visual layouts (background templates) present in the training distribution, explaining the validation-test generalization gap.

\subsubsection{Qualitative Rationale Evaluation (VLM Human-in-the-Loop Rubric)}
For generative VLMs (Qwen2-VL-7B and Gemini 1.5 Pro), explainability is evaluated through natural-language justifications. We conduct a human-in-the-loop pilot evaluation on 30 representative memes from the validation split. Explanations are rated from 1 (Poor) to 5 (Excellent) across three dimensions:
\begin{itemize}
    \item \textbf{Correctness} (Average: 4.1 / 5.0): Evaluates whether the generated rationale is logically consistent with the ground truth label. Deductions are primarily driven by false negatives.
    \item \textbf{Specificity} (Average: 4.8 / 5.0): Evaluates how well the model references specific text overlays or visual objects. Anchoring the VLM prompt with OCR transcripts prevented text misreading, yielding very high scores.
    \item \textbf{Usefulness} (Average: 4.5 / 5.0): Evaluates whether the explanation provides actionable, readable summaries that would help a human content moderator quickly make a flag decision.
\end{itemize}

\subsection{The Role of Closed APIs as Research References}
We report Gemini 1.5 Pro zero-shot as a reference baseline. It serves as an upper bound, achieving 0.8829 on Task A without local training, validating the general-reasoning capabilities of commercial foundation models.

\section{Conclusion}
\label{sec:conclusion}
We presented a systematic benchmark of over 30 models on MAMI 2022. While Gemini 1.5 Pro zero-shot sets SOTA on Task A, a simple SVM-RBF on frozen CLIP features matches or beats SOTA on Task B. Crucially, our SHAP modality analysis (72.42\% visual reliance) explains the validation-test generalization gap by showing that trained classifiers rely on template memorization. Future work should focus on debiased, template-aware splits.

\bibliographystyle{IEEEtran}
\bibliography{references}

\end{document}
